\documentclass[10pt,a4paper]{amsart}
\usepackage{amsmath,amssymb,amsthm,graphicx}
% Retained minimal preamble for the Stacks Project chapter more-algebra.tex (More on Algebra).
% The chapter file pinned-original-source.tex carries only \input{preamble} on its first line, so no
% preamble text lives inside the chapter itself.  The Stacks Project uses one single shared preamble for
% every chapter; the snapshot stages it identically as sources/stacks-chow/preamble.tex and
% sources/stacks-etale-cohomology/preamble.tex.  The block below is copied mechanically, line for line,
% from sources/stacks-chow/preamble.tex lines 14-18 (the Xy-pic setup) and 154-217 (the native theorem
% declarations together with the mathematical macros).  Nothing was added, removed or redefined; the
% wrapper class and the amsmath/amssymb/amsthm/graphicx base packages are supplied by the assembler.
% Chapter: more-algebra.tex -- More on Algebra.
% For commutative diagrams we use Xy-pic
\usepackage[all]{xy}

% We use 2cell for 2-commutative diagrams.
\xyoption{2cell}

% Theorem environments.
%
\theoremstyle{plain}
\newtheorem{theorem}[subsection]{Theorem}
\newtheorem{proposition}[subsection]{Proposition}
\newtheorem{lemma}[subsection]{Lemma}

\theoremstyle{definition}
\newtheorem{definition}[subsection]{Definition}
\newtheorem{example}[subsection]{Example}
\newtheorem{exercise}[subsection]{Exercise}
\newtheorem{situation}[subsection]{Situation}

\theoremstyle{remark}
\newtheorem{remark}[subsection]{Remark}
\newtheorem{remarks}[subsection]{Remarks}

\numberwithin{equation}{subsection}

% Macros
%
\def\lim{\mathop{\mathrm{lim}}\nolimits}
\def\colim{\mathop{\mathrm{colim}}\nolimits}
\def\Spec{\mathop{\mathrm{Spec}}}
\def\Hom{\mathop{\mathrm{Hom}}\nolimits}
\def\Ext{\mathop{\mathrm{Ext}}\nolimits}
\def\SheafHom{\mathop{\mathcal{H}\!\mathit{om}}\nolimits}
\def\SheafExt{\mathop{\mathcal{E}\!\mathit{xt}}\nolimits}
\def\Sch{\mathit{Sch}}
\def\Mor{\mathop{\mathrm{Mor}}\nolimits}
\def\Ob{\mathop{\mathrm{Ob}}\nolimits}
\def\Sh{\mathop{\mathit{Sh}}\nolimits}
\def\NL{\mathop{N\!L}\nolimits}
\def\CH{\mathop{\mathrm{CH}}\nolimits}
\def\proetale{{pro\text{-}\acute{e}tale}}
\def\etale{{\acute{e}tale}}
\def\QCoh{\mathit{QCoh}}
\def\Ker{\mathop{\mathrm{Ker}}}
\def\Im{\mathop{\mathrm{Im}}}
\def\Coker{\mathop{\mathrm{Coker}}}
\def\Coim{\mathop{\mathrm{Coim}}}

% Boxtimes
%
\DeclareMathSymbol{\boxtimes}{\mathbin}{AMSa}{"02}

%
% Macros for moduli stacks/spaces
%
\def\QCohstack{\mathcal{QC}\!\mathit{oh}}
\def\Cohstack{\mathcal{C}\!\mathit{oh}}
\def\Spacesstack{\mathcal{S}\!\mathit{paces}}
\def\Quotfunctor{\mathrm{Quot}}
\def\Hilbfunctor{\mathrm{Hilb}}
\def\Curvesstack{\mathcal{C}\!\mathit{urves}}
\def\Polarizedstack{\mathcal{P}\!\mathit{olarized}}
\def\Complexesstack{\mathcal{C}\!\mathit{omplexes}}
% \Pic is the operator that assigns to X its picard group, usage \Pic(X)
% \Picardstack_{X/B} denotes the Picard stack of X over B
% \Picardfunctor_{X/B} denotes the Picard functor of X over B
\def\Pic{\mathop{\mathrm{Pic}}\nolimits}
\def\Picardstack{\mathcal{P}\!\mathit{ic}}
\def\Picardfunctor{\mathrm{Pic}}
\def\Deformationcategory{\mathcal{D}\!\mathit{ef}}
\usepackage{cleveref}
\crefname{originallemma}{original lemma}{original lemmas}
\Crefname{originallemma}{Original lemma}{Original lemmas}
\makeatletter
\newlabel{more-algebra-lemma-base-change-H1-regular}{{063Q}{source}{}{sourcepointer-1}{}}
\newlabel{more-algebra-lemma-base-change-H1-regular@cref}{{[originallemma][0][]063Q}{[1][0][]source}{}{}{}}
\makeatother
\crefname{originallemma}{original lemma}{original lemmas}
\Crefname{originallemma}{Original lemma}{Original lemmas}
\makeatletter
\newlabel{algebra-lemma-flat-tor-zero}{{00HL}{source}{}{sourcepointer-2}{}}
\newlabel{algebra-lemma-flat-tor-zero@cref}{{[originallemma][0][]00HL}{[1][0][]source}{}{}{}}
\makeatother
\title{The first Koszul homology of a sequence with flat quotient maps onto its base change}
\author{The Stacks Project Authors}
\date{stacks-project:more-algebra.tex:local-snapshot-20261012; GNU FDL 1.2 or later}
\begin{document}
\maketitle
\setcounter{subsection}{1}
\noindent\textit{Editorial note.} All delivered mathematical text is the author's own, copied from the pinned Stacks Project chapter file more-algebra.tex; no statement or proof step was rewritten, reordered or removed.\par\medskip
\noindent\textit{Editorial note.} The chapter file pinned-original-source.tex carries only \textbackslash{}input\{preamble\} on its first line, so no preamble text lives inside the chapter. The retained file sources/stacks-more-algebra/retained-preamble.tex is a mechanical, line-for-line copy of the project's single shared preamble as staged in the snapshot, namely sources/stacks-chow/preamble.tex lines 14-18 (Xy-pic setup) and 154-217 (theorem declarations and mathematical macros). Nothing was added, removed or redefined, no package was added, and the excerpt serves only as preamble, never as a span coordinate.\par\medskip
\noindent\textit{Editorial note.} Both counted spans start at the \textbackslash{}begin line of their environment, so the statement and the author's proof are delivered whole; surrounding prose of the section is deliberately outside the counted spans and is carried as standard prerequisites instead.\par\medskip
\noindent\textit{Editorial note.} The citation fields (source\_id, source\_version, authors, source\_url, license.name, source\_title) are identical to those of the cards already published from this chapter, so that the batch selector treats them as one source instead of deferring them as a same-source clash.\par\medskip
\noindent\textit{Allowed standard prerequisites (editorial).}\begin{itemize}
\item Koszul complex of a sequence and its first homology, together with the four-term presentation ending in B/J. used at proof lines 7949-7964, where the complex and the three short exact sequences it is split into are written down. Provides the objects whose behaviour under tensor with A' is being tracked.
\item If B/J is flat over A then tensoring with A' preserves the exactness of 0 to J to B to B/J to 0, i.e. vanishing of the relevant Tor. the author cites this at proof lines 7974-7975 as Lemma algebra-lemma-flat-tor-zero of the Algebra chapter (tag 00HL). Is the single hypothesis that makes the first tensored sequence exact, which in turn gives J' = J tensor A' and the rest of the computation.
\item Right exactness of the tensor product and the identification of the first Koszul homology with the cokernel of the map to the kernel. used at proof lines 7964-7973 and 7979-7987. Turns the exactness information into the surjectivity of the canonical map H\_1 tensor A' to H'\_1 asserted by the statement.
\end{itemize}\medskip
\ifdefined\hypertarget\hypertarget{sourcepointer-1}{}\fi
\noindent\textit{Source reference.} Original lemma 063Q: focused-seeds/AGP/sources/stacks-snapshot-20261011/tags.txt, source lines 7839--7839. The lemma's own canonical tag, quoted so that readers can resolve the label inside the delivered text; it points at the counted statement itself and carries no proof content.\par\medskip
\ifdefined\hypertarget\hypertarget{sourcepointer-2}{}\fi
\noindent\textit{Source reference.} Original lemma 00HL: focused-seeds/AGP/sources/stacks-snapshot-20261011/tags.txt, source lines 647--647. Navigation pointer into the Algebra chapter for the fact that a flat quotient makes the relevant Tor vanish; it is the standard prerequisite that lets the first tensored sequence stay exact, and every subsequent step of the computation is inside the counted proof span.\par\medskip
\begin{lemma}
\label{lemma-base-change-H1-regular}
Let $A \to B$ be a ring map.
Let $f_1, \ldots, f_r$ be a sequence in $B$ such that $B/(f_1, \ldots, f_r)$
is $A$-flat. Let $A \to A'$ be a ring map. Then the canonical map
$$
H_1(K_\bullet(B, f_1, \ldots, f_r)) \otimes_A A'
\longrightarrow
H_1(K_\bullet(B', f'_1, \ldots, f'_r))
$$
is surjective. Here $B' = B \otimes_A A'$ and $f_i' \in B'$ is the image
of $f_i$.
\end{lemma}

\begin{proof}
The sequence
$$
\wedge^2(B^{\oplus r}) \to B^{\oplus r} \to B \to B/J \to 0
$$
is a complex of $A$-modules with $B/J$ flat over $A$ and
cohomology group $H_1 = H_1(K_\bullet(B, f_1, \ldots, f_r))$ in the spot
$B^{\oplus r}$. If we tensor this with $A'$ we obtain a complex
$$
\wedge^2((B')^{\oplus r}) \to (B')^{\oplus r} \to B' \to B'/J' \to 0
$$
which is exact at $B'$ and $B'/J'$. In order to compute its
cohomology group $H'_1 = H_1(K_\bullet(B', f'_1, \ldots, f'_r))$
at $(B')^{\oplus r}$ we split the first sequence above into
the exact sequences $0 \to J \to B \to B/J \to 0$,
$0 \to K \to B^{\oplus r} \to J \to 0$, and
$\wedge^2(B^{\oplus r}) \to K \to H_1 \to 0$.
Tensoring over $A$ with $A'$ we obtain the exact sequences
$$
\begin{matrix}
0 \to J \otimes_A A' \to B \otimes_A A' \to (B/J) \otimes_A A' \to 0 \\
K \otimes_A A' \to B^{\oplus r} \otimes_A A' \to J \otimes_A A' \to 0 \\
\wedge^2(B^{\oplus r}) \otimes_A A' \to K \otimes_A A' \to H_1 \otimes_A A'
\to 0
\end{matrix}
$$
where the first one is exact as $B/J$ is flat over $A$, see
Algebra, Lemma \ref{algebra-lemma-flat-tor-zero}. We conclude
that $J' = J \otimes_A A' \subset B'$ and that
$K \otimes_A A' \to \Ker((B')^{\oplus r} \to B')$ is
surjective. Thus
\begin{align*}
H_1 \otimes_A A'
& =
\Coker\left(\wedge^2(B^{\oplus r}) \otimes_A A' \to K \otimes_A A'\right) \\
& \to
\Coker\left(
\wedge^2((B')^{\oplus r})  \to \Ker((B')^{\oplus r} \to B')
\right) = H'_1
\end{align*}
is surjective too.
\end{proof}

\end{document}
